\documentclass[11pt,a4paper]{article}
\usepackage[margin=2.2cm]{geometry}
\usepackage{amsmath,amssymb,xcolor,enumitem}
\usepackage[T1]{fontenc}\usepackage{mathptmx}
\pdfminorversion=4
\newcommand{\rev}[1]{\par\smallskip\noindent\textbf{Reviewer:} \emph{#1}\par\smallskip\noindent\textbf{Response:} }
\begin{document}
\begin{center}
{\Large\bfseries Response to the Reviewer}\\[4pt]
Flatness and Passivity-Based Control with Adaptive Neural Compensation of a Direct-Drive PMSG Wind Energy Conversion System\\[2pt]
S. Regani, M. Messadi, K. Kemih
\end{center}

We thank the reviewer for the careful reading and the constructive comments. All of them have been
addressed. Section, equation, table and figure numbers refer to the revised manuscript.

\section*{Major comments}

\rev{1. Originality: the paper combines known techniques; the contributions must be stated clearly.}
We agree that flatness, PBC and adaptive neural compensation are individually known. The Introduction
now ends with an explicit list of contributions: (i) a DC-link flat output that includes the energy of
the $L$ filter, which gives an exact generator-power feedforward and linear error dynamics;
(ii) PBC current laws for both converters and a cascade proof that the complete nominal interconnection
is exponentially stable (new Proposition~2); (iii) the transfer of the neural compensation of
our earlier work from the rotor currents of a DFIG to the two flat loops of a PMSG;
(iv) a quantitative comparison with PI vector control and an extended validation.

\rev{2. Sign inconsistency in the DC-link control derivation.}
We thank the reviewer for this remark. The implemented law was correct, but the text stated the
imposed error dynamics with the wrong sign. The derivation in Section~V-A now states that the law
imposes $\dot e_y=-\nu_2$, i.e. $\dot y_2=\nu_2$, and shows explicitly that inserting the law in the
$y_2$ dynamics gives $\dot y_2=\nu_2-\Delta P_g$, hence
$\dot e_y+k_{21}e_y+k_{22}\int e_y=\Delta P_g$. We also state that $\dot y_2^*$ is neglected because
$V_{dc}^*$ is constant and the filter energy varies slowly. The signs of all equations were checked
against the simulation code.

\rev{3. Stability of the whole interconnection.}
Proposition~1 (PBC, Barbalat) only shows convergence of the current errors. We added Proposition~2 with
its proof: with $V_s=H+\varepsilon\,\mathbf e^T L\mathbf z$ the inner loops are exponentially stable
(for bounded electrical speed), and the outer flat loops are exponentially stable linear systems driven
by $\frac{3p\psi_f}{2J}e_{sq}$ and $\frac32 v_{gd}e_{gd}$. The cascade of exponentially stable
systems is exponentially stable. This provides, locally (outside the saturations), the assumption
required by the uniform ultimate boundedness result of the neural compensation. The assumptions
(exact model, no saturation, ideal PLL) are stated explicitly.

\rev{4. Comparison with conventional (PI) and/or nonlinear controllers.}
A new Section~VI compares the proposed control with PI vector control (PI speed, DC-voltage and current
loops with $dq$ decoupling) tuned on the same poles, with and without power feedforward (PI+FF), in the
nominal and robustness cases (Table~III, Fig.~3). Under parameter errors the proposed control tracks
the speed 3.8 times more accurately (IAE 0.0109/0.0068 rad vs 0.0419 rad) and limits the DC-link
deviation during a 20\,\% grid dip to 3.6~V vs 11.6~V (PI+FF) and 62~V (PI). The table also reports
the drawbacks honestly: PI gives smaller current errors, and PI+FF a smaller start-up DC-link
deviation. The PI tuning and the common MPPT reference and limits are described in Section~VI-A.

\rev{5. More extensive validation (wind profiles, initial conditions, switched model).}
Added: five turbulence realizations (Table~IV, mean and spread), an initial speed 10\,\% below the
MPPT value (Table~IV) and 10\,\% above it (text), and a switched model with SVPWM at 2160~Hz and a
1~$\mu$s step at 8, 9 and 10~m/s (Section~VI-D) with the grid-current THD (4.4, 3.1 and 2.3\,\%,
below the 5\,\% limit of IEEE~519).

\rev{6. ANN trade-offs and sensitivity to $\Gamma$, $\delta$, $\sigma$.}
Table~IV gives the sensitivity to $\Gamma\in\{1,2,5,10,20\}$, $\delta\in\{0.05,0.1,0.2\}$ and
$\sigma\in\{0.001,0.01,0.1\}$. A larger $\Gamma$ improves the speed tracking but increases the current
errors and becomes oscillatory at $\Gamma=20$; $\Gamma=5$ is retained as a compromise. During this study
we found that the original constant-gain update could oscillate for large initial errors; the update
laws now use the normalized gain $\Gamma/(1+\|\mathbf z\|^2)$, which removes the problem and keeps the
boundedness proof. The text states clearly that the network mainly improves the speed loop
($-38$\,\% IAE) and slightly increases the current errors and the dip deviation.

\section*{Minor comments}
\begin{enumerate}[leftmargin=*]
\item \emph{Readability of Fig.~2.} The block diagram was redrawn with larger fonts and a clearer
signal flow.
\item \emph{Typography of Eq.~(22).} Corrected (rewritten as an aligned equation).
\item \emph{Normalization of the ANN inputs.} Section~V-C now gives all scaling factors (speed,
energy and current errors, filter time constant of the derivatives and their scaling).
\item \emph{Status of reference [15].} Updated to its current status (submitted to \emph{Mathematics}; it will be replaced by the full citation if accepted before the camera-ready version).
\item \emph{THD computation.} Specified: FFT of the phase-$a$ grid current over ten grid periods in
steady state, harmonics up to the 40th order.
\item \emph{Computational complexity.} Section~VI-D now gives the operation count (about 200
multiply--accumulate operations and ten $\tanh$ evaluations per step, i.e. a few MFLOP/s at 10~kHz),
compatible with a standard DSP.
\end{enumerate}
\end{document}
